\begin{figure}[t]
\centering
\begin{tikzpicture}[
    >=Latex,
    font=\small,
    axis/.style={->,thick},
    domain/.style={draw=blue!65!black,thick,fill=blue!5},
    cell/.style={draw=black!45,thin},
    hilite/.style={draw=green!50!black,thick,fill=green!18},
    note/.style={rounded corners,draw=black!40,fill=white,inner sep=4pt,align=center}
]

% ------------------------------------------------------------------
% Left panel: continuum measure
% ------------------------------------------------------------------
\node[font=\bfseries] at (2.0,4.65) {Continuum measure};

\draw[axis] (0,0) -- (4.45,0) node[right] {$u_A$};
\draw[axis] (0,0) -- (0,4.45) node[above] {$u_B$};
\draw[domain] (0,0) -- (0,4) -- (4,0) -- cycle;
\node[anchor=east] at (-0.08,4) {$U$};
\node[anchor=north] at (4,-0.08) {$U$};
\node[anchor=north east] at (-0.05,-0.05) {$0$};

% A schematic non-atomic measure cloud
\fill[blue!35,opacity=.28] (1.15,1.15) circle (0.72);
\fill[cyan!55,opacity=.24] (1.55,0.82) circle (0.58);
\fill[orange!65,opacity=.20] (0.85,1.55) circle (0.48);
\fill[blue!55,opacity=.18] (2.35,0.72) circle (0.42);
\node[align=center] at (1.25,2.30)
    {$\nu$ arbitrary on\\[-1mm] $D_U=\{u_A+u_B\le U\}$};

\node[note,text width=4.15cm] at (2.05,-1.05)
    {$\displaystyle
      \Lambda_{st}
      =\int_{D_U} K_{\tau}^{st}(u)\,d\nu(u)
      +r^{\rm tail}_{st}$\\
     $\displaystyle |r^{\rm tail}_{st}|\le\varepsilon$};

% Transition arrow
\draw[->,very thick,gray!65] (4.85,2.0) -- (6.10,2.0);
\node[align=center,text=black!70] at (5.48,2.55)
    {\footnotesize partition\\[-1mm]\footnotesize into cells};

% ------------------------------------------------------------------
% Right panel: cell masses
% ------------------------------------------------------------------
\begin{scope}[xshift=6.55cm]
\node[font=\bfseries] at (2.0,4.65) {Finite cell variables};

\draw[axis] (0,0) -- (4.45,0) node[right] {$u_A$};
\draw[axis] (0,0) -- (0,4.45) node[above] {$u_B$};
\draw[domain] (0,0) -- (0,4) -- (4,0) -- cycle;
\node[anchor=east] at (-0.08,4) {$U$};
\node[anchor=north] at (4,-0.08) {$U$};
\node[anchor=north east] at (-0.05,-0.05) {$0$};

% Grid, clipped to triangular domain
\begin{scope}
\clip (0,0) -- (0,4) -- (4,0) -- cycle;
\foreach \x in {0.8,1.6,2.4,3.2}
    \draw[cell] (\x,0) -- (\x,4);
\foreach \y in {0.8,1.6,2.4,3.2}
    \draw[cell] (0,\y) -- (4,\y);
\end{scope}

% Highlight one complete cell
\filldraw[hilite] (0.8,0.8) rectangle (1.6,1.6);
\node at (1.04,1.40) {$C_j$};
\fill[black] (1.20,1.20) circle (1.5pt);
\node[anchor=west] at (1.27,1.19) {$c_j$};

% Show that the actual mass need not sit at c_j
\fill[green!45!black] (0.94,0.98) circle (1.1pt);
\fill[green!45!black] (1.48,1.08) circle (1.1pt);
\fill[green!45!black] (1.36,1.48) circle (1.1pt);

\node[note,text width=3.15cm] (massnote) at (2.85,2.85)
    {$q_j:=\nu(C_j)$\\[1mm]
     mass may lie anywhere in $C_j$};
\draw[->,green!45!black,thick]
    (massnote.south west) .. controls (2.25,2.25) and (1.85,1.75) .. (1.53,1.47);

\node[note,text width=4.35cm] at (2.05,-1.05)
    {$\displaystyle
      \int_{C_j}K_{\tau}^{st}(u)\,d\nu(u)
      =q_jK_{\tau}^{st}(c_j)+r_{st,j}$\\
     $\displaystyle |r_{st,j}|\le q_j\delta_{st,j}$};
\end{scope}

% ------------------------------------------------------------------
% Certified global constraint
% ------------------------------------------------------------------
\node[
    rounded corners,
    draw=orange!65!black,
    fill=orange!6,
    inner sep=7pt,
    text width=11.2cm,
    align=center
] at (5.30,-2.55)
{
\[
\left|
\sum_z G[(z,s),(z,t)]
-
\sum_j q_j K_{\tau}^{st}(c_j)
\right|
\le
\sum_j q_j\delta_{st,j}
+\varepsilon ,
\qquad
\delta_{st,j}\ge
\sup_{u\in C_j}
\left|K_{\tau}^{st}(u)-K_{\tau}^{st}(c_j)\right|.
\]
};

\node[
    rounded corners,
    draw=green!45!black,
    fill=green!5,
    inner sep=5pt,
    text width=10.7cm,
    align=center
] at (5.30,-3.95)
{
Every continuum measure maps to feasible cell variables:
\[
q_j=\nu(C_j),
\qquad
\varepsilon=\nu(D_U^c),
\qquad
\mathcal F_{\rm continuum}\subseteq\mathcal F_{\rm cell}.
\]
The representative points \(c_j\) are only used to evaluate the kernel;
they are \emph{not} imposed support points of \(\nu\).
};

\end{tikzpicture}
\caption{Schematic of the certified continuum-to-cell relaxation in Stage~3.
The continuum measure is replaced by cell masses \(q_j=\nu(C_j)\), not by
point masses at the representatives \(c_j\).  The variables
\(\delta_{st,j}\) bound the possible kernel variation within each cell, while
\(\varepsilon\) bounds the omitted tail.  Consequently every feasible
continuum measure produces a feasible point of the finite conic program,
which is the required outer-approximation direction.  The figure is written
with the exact finite-\(\tau\) kernel \(K_\tau\); the same construction applies
to the limiting kernel \(K_0\).}
\label{fig:stage3-cell-relaxation}
\end{figure}
